
\subsubsection{6.1 Summary of Findings}

\textbf{Finding 1: Behavioral proxy signals are detectable in returning-user cohorts, but in the direction opposite to the BAA disengagement prediction.} The positive prompt length trend ($\tau$ = +0.744) establishes that behavioral signals are computationally accessible in public interaction logs --- but contradicts the BAA prediction of $\tau$ < 0.

\textbf{Finding 2: Two of three proxies are unmeasurable under current conditions.} Vote entropy is blocked by dataset access restrictions; correction frequency has an insufficient signal base rate. Only one proxy is evaluable, and it shows a direction opposite to the BAA prediction.

\textbf{Finding 3: The returning-user cohort design introduces a selection bias that prevents causal attribution.} The $\geq 3$ monthly bin filter is necessary for signal detection but systematically retains high-engagement users, preventing generalization to the broader user population.

\subsubsection{6.2 Why Is Prompt Length Increasing?}

Three competing explanations are consistent with the positive trend, all independent of any AI quality improvement mechanism. The data available in this study cannot discriminate among them.

\textbf{Explanation 1: Returning-user selection bias (most plausible).} High-engagement users who return $\geq 3$ months are power users who compose longer prompts as a baseline behavioral profile. The positive trend may reflect increasing dominance of heavy users over time rather than within-user behavioral change. \textit{Discriminating test:} compare prompt length trajectory for users appearing in their first month versus users in their third or later month of activity in a difference-in-differences design.

\textbf{Explanation 2: User expertise gain (also plausible).} Users learn to compose more effective prompts --- longer, more context-rich --- over time. This expertise gain predicts increasing prompt length independent of AI quality. \textit{Discriminating test:} within-user prompt length trajectory across monthly appearances (requires individual-level longitudinal data not recoverable from IP-hash pseudonymization).

\textbf{Explanation 3: Platform adoption effect (moderate plausibility).} WildChat's user base may have shifted toward professional or developer use during the observation period, involving longer technical prompts. \textit{Discriminating test:} topic-tag stratification of the prompt length trend using the topic tags available in WildChat-1M.

The BAA disengagement mechanism requires decreasing prompt length. All three explanations above predict increasing prompt length for reasons unrelated to AI quality. The study cannot rule out any of these; the comparison group experiment described in Explanation 1 is the most direct test.

\subsubsection{6.3 Limitations}

\textbf{L1: Returning-user selection bias (severity: HIGH).} The analysis cohort (n = 27,902) is self-selected for high engagement by construction of the $\geq 3$ monthly bin filter. This is not representative of the broader WildChat user population. The proposed control experiment is a two-group comparison: returning users ($\geq 3$ bins) versus first-appearance users, matched by month and topic, in a difference-in-differences design.

\textbf{L2: LMSYS primary dataset access (severity: HIGH for P2).} Vote entropy --- arguably the most direct operationalization of BAA preference discrimination behavior --- is entirely blocked by dataset access restrictions. Remediation requires requesting approved research access to \texttt{lmsys/chatbot\textbackslash \_arena\textbackslash \_conversations}.

\textbf{L3: Correction frequency proxy inadequacy (severity: MEDIUM).} Regex pattern matching captures explicit corrections in approximately 0.05\% of turns. Implicit corrections (rephrasing, session abandonment, resubmission) require LLM-based annotation or session-level behavioral proxies. The current proxy is underpowered regardless of the true underlying trend.

\textbf{L4: No causal chain verification (severity: HIGH).} The BAA mechanism posits three steps: (1) AI quality improves, (2) users perceive less need to probe carefully, (3) behavioral engagement declines. None of these steps is verified. Most critically, WildChat contains no per-interaction measure of AI model quality or version; the causal link from AI improvement to observed behavioral trends cannot be established. Future work must include an AI quality covariate (e.g., per-bin model version or quality rating) to test the causal chain.

\textbf{L5: IP-hash cohort noise (severity: MEDIUM).} Multiple users may share a single IP address (NAT, shared networks); a single user may appear under multiple IP addresses (VPN, mobile IP reassignment). The cohort (n = 27,902) is defined by IP-hash continuity. The magnitude of this noise is not quantified in the present study.

\textbf{L6: WildChat platform scope (severity: LOW--MEDIUM).} WildChat logs capture ChatGPT (OpenAI API) interactions proxied through the WildChat web interface. Results may not generalize to other AI platforms (e.g., Claude, Gemini) or direct API users.

\textbf{L7: GPT-3.5 $\rightarrow$ GPT-4 model transition confound (severity: MEDIUM).} GPT-4 API availability expanded in March 2023 --- immediately before the observation window \citep{April2023}. Users transitioning from GPT-3.5-turbo to GPT-4 would encounter substantially larger context windows and improved instruction following, providing a technical incentive to compose longer and more detailed prompts regardless of engagement level. This confound partially overlaps with L4 but is more specific: the GPT-3.5 $\rightarrow$ GPT-4 transition is a discrete event whose timing aligns with the start of the observation window. Disentangling model capability effects from user behavioral adaptation requires stratifying by model version within WildChat --- a stratification not performed in the present study.

\subsubsection{6.4 Broader Implications}

This work develops and validates methods for longitudinal behavioral analysis in large-scale AI interaction logs. The returning-user cohort construction pipeline, Hamed-Rao Mann-Kendall testing, and bootstrap confidence interval framework are applicable to any dataset with user-level pseudonymization and sufficient temporal coverage. The analysis operates exclusively at the cohort aggregate level; the IP-hash approach does not enable individual re-identification.

The high autocorrelation in the prompt token count series (ACF lag-1 = 0.634) is itself informative: it indicates that monthly cohort-level behavioral signals are persistent rather than noisy, suggesting that behavioral dynamics in returning-user populations have detectable temporal structure. Standard Mann-Kendall without Hamed-Rao correction would yield a materially different (overestimated) significance assessment in this setting.

